%*****************************************
\chapter{Ergebnisse}\label{ch:results}
%*****************************************

Im Folgenden werden die auf Grundlage der durchgeführten Interviews und der bisherigen Überlegungen erstellten Modelle und Diagramme vorgestellt. Dabei werden zunächst die einzelnen Interessengruppen aufgelistet, die in dem Umfeld der medizinischen Versorgung von Wohnungslosen Einfluss ausüben. Mithilfe von Kommunikationsdiagrammen der \ac{UML} werden einige konkrete Angebote oder Gegebenheiten näher beleuchtet und anhand des 3-Ebenen-Metamodells (\acs{3LGM²}) wird ein Gesamtbild des Leipziger Informationssystems für die medizinische Dokumentation bei dieser Personengruppe beschrieben.

\section{Stakeholder-Analyse}\label{sec:stakeholder}

In Vorbereitung auf einen Vergleich und zusammenhängend mit den zuvor definierten Kriterien werden zunächst die Interessen näher untersucht, die von den einzelnen Parteien an die medizinische Versorgung von Wohnungslosen gestellt werden. Die Anforderungen an ein Dokumentationssystem, welches hierbei eine unterstützende Funktion besitzt, sollten im Einklang mit diesen Punkten gewählt werden. Die durchgeführte Stakeholder-Analyse ist in Tabelle \ref{tab:stakeholder} zu finden.

\begin{table}[ht]
	\centering
	\begin{tabulary}{\textwidth}{Lp{0.7\linewidth}}
		\toprule
		Stakeholder&				Interessen\\
		\midrule
		Betroffene&					einfacher Zugang zu medizinischer Versorgung\\
		&							Vermeidung von Fehl- oder Mehrfachbehandlungen / -untersuchungen\\
		&							ggf. Wahrung der Anonymität\\
		Einwohner&					keine Beeinträchtigung ihrer med. Versorgung\\
		CABL und UVO&				Verbesserung der med. Versorgung von Wohnungslosen\\
		&							möglichst vollständige Dokumentation zu Fallzahlen u.ä.\\
		Safe&						Verbesserung der Lebenssituation von Wohnungslosen\\
		Sozialarbeiter:innen&		Beibehaltung der persönlichen Kontakte\\
		&							Belastungsreduzierung\\
		Ärzt:innen&					verfügbare Dokumentation für Diagnose und Behandlung\\
		&							Belastungsreduzierung / kein wesentlicher Mehraufwand\\
		Sozialamt&					Erfüllung der Aufgaben zur Sozialhilfe\\
		&							Kostenminimierung\\
		Gesundheitsamt&				Erfüllung der Aufgaben zur Gesundheitsverwaltung\\
		&							Kostenminimierung\\
		\bottomrule
	\end{tabulary}
	\caption[Stakeholder]{Übersicht der Stakeholder}
	\label{tab:stakeholder}
\end{table}

Es ist anzumerken, dass sich die Gruppe der Betroffenen auf diejenigen Personen bezieht, die nach aktuellem Stand die Angebote der verschiedenen Hilfsorganisationen in Anspruch nehmen. Sie besteht dahingehend nicht nur aus Wohnungs- und Obdachlosen, sondern beinhaltet auch EU-Bürger oder Drittstaatler, die beispielsweise bei der regulären medizinischen Versorgung vernachlässigt werden, wenn keine Anbindung an eine Krankenversicherung möglich ist. Das Hauptinteresse ist deswegen der Zugang zu medizinischen Angeboten, die qualitativ möglichst identisch mit denen der Regelversorgung ist, wobei es einige Personen vermeiden wollen, ihre Identität preiszugeben.

Für den allgemeinen Bürger bzw. Einwohner, welcher über eine Krankenversicherung abgesichert ist, ist es wichtig, dass bestehende Möglichkeiten für eine medizinische Versorgung nicht eingegrenzt werden und die terminliche und örtliche Verfügbarkeit weiterhin gewährleistet wird.

Die verschiedenen Hilfsorganisationen und Projekte verfolgen nach wie vor die Ziele, die (medizinische) Versorgung bzw. Lebenssituation der Betroffenen zu verbessern und würden auf Änderungen, die diese angehen, positiv gegenüberstehen. Es wäre in diesem Zusammenhang auch wünschenswert, dass bestimmte Grunddaten wie z.B. Fallzahlen mit erfasst werden können, um statistische Auswertungen durchzuführen.

Bei sozialarbeitendem und medizinischem Personal wird eine Reduzierung der Belastung angestrebt, was insbesondere von Bedeutung ist, wenn es sich um ehrenamtliche Tätigkeiten handelt, da dabei der entstehende Mehraufwand möglichst verringert werden sollte. Angesichts ihrer menschennahen Tätigkeit, ist es für Sozialarbeiter:innen bzw. explizit Streetworker:innen von Bedeutung, dass die aufgebauten Kontakte zu Wohnungslosen bestehen bleiben.

Das Sozial- sowie das Gesundheitsamt streben zunächst eine Erfüllung der ihnen aufgetragenen Aufgaben an, wobei als sekundärer Betrachtungspunkt eine Minimierung der anfallenden Kosten hinzugezogen wird. Die Tätigkeitsschwerpunkte sind dabei breit gefächert, beispielsweise setzt sich das Aufgabenspektrum des Gesundheitsamtes aus Verpflichtungen gegenüber der Bevölkerungs- und Individualmedizin zusammen. Mit anderen Worten erfolgt nicht nur eine Einschätzung und Überwachung der allgemeinen gesundheitlichen Situation der Bevölkerung, sondern auch die Koordination von verschiedenen medizinischen Hilfsangeboten.


\section{Kommunikationsmodelle}\label{sec:com}

Die nachfolgenden Kommunikationsmodelle stellen den Informationsfluss in den ausgewählten Angeboten dar, die Wohnungslosen und Menschen, die nicht über eine Krankenkasse abgesichert sind, als Möglichkeit einer medizinischen Versorgung zur Verfügung stehen. Da diese Art von Diagramm in enger Verbindung mit dem Sequenzdiagramm steht, wird hierbei auch ein zeitlicher Ablauf mit dargestellt. Es ist dennoch anzumerken, dass primär die Informationsflüsse und die übermittelten Daten im Vordergrund stehen und dadurch im Fokus der Erklärungen stehen.

\subsection{CABL-Sprechstunden}\label{sub:cabl}

\begin{figure}[h]
	\centering
	\includegraphics[width=\linewidth]{Images/cd_cabl}
	\caption[Kommunikationsdiagramm CABL-Sprechstunde]{Kommunikationsdiagramm CABL-Sprechstunde}
	\label{fig:cdCABL}
\end{figure}

Die Sprechstunden von \ac{CABL} gehören zu dem Versuch, hilfesuchende Personen unabhängig ihrer Umstände den Zugang zu Gesundheitsdiensten zu ermöglichen. Bei dem Clearing sind neben den jeweiligen Betroffenen auch ein Sozialarbeiter oder eine Sozialarbeiterin des Vereins sowie jemand aus dem ärztlichen Ehrenamt mit anwesend. Der Verein \ac{AKST} stellt dafür das Anwendungssystem \enquote{Schein-Verwaltung} zur Verfügung.

Wenn eine Person die Sprechstunde aufsucht, so wird zum einen versucht, ein zuständiger Kostenträger zu ermitteln oder anderweitig eine Absicherung über eine Krankenkasse zu erreichen. Zum anderen wird mittels Anamnese die Fachrichtung des jeweiligen Falls ermittelt. Die Daten, die während der Sprechstunde anfallen, werden in die Schein-Verwaltung eingetragen sowie mittels verschiedener einfacher Tools wie Excel-Listen oder Fragebögen dokumentiert.

Kann kein Kostenträger ausfindig gemacht werden oder ist das Anliegen von Dringlichkeit, so wird das Clearing abgeschlossen, ein \ac{ABS} erstellt und ausgehändigt. Die betroffene Person kann anschließend bei einer Gesundheitsinstitution mit dem Fachbereich, für den der \ac{ABS} ausgestellt wurde, vorstellig werden, wobei sich das Verfahren ab diesen Punkt mit dem eines Krankenkassenpatienten gleicht. Zuletzt kann der \ac{AKST}, da dieser das Tool dafür bereitstellt, auch eine Auswertung der Schein-Verwaltungs-Daten durchführen.

\subsection{UVO-Koordination}\label{sub:uvo}

\begin{figure}[h]
	\centering
	\includegraphics[width=\linewidth]{Images/cd_uvo}
	\caption[Kommunikationsdiagramm UVO-Angebote]{Kommunikationsdiagramm zu den UVO-Angeboten}
	\label{fig:cdUVO}
\end{figure}

Das Projekt \ac{UVO}, welches dem Verein \ac{CABL} unterstellt ist, bietet wie bereits erwähnt aufsuchende Angebote und feste Sprechstunden mit ehrenamtlichen Ärzt:innen an. Um die Inanspruchnahme und die allgemeine Wirkung dieser Aktionen adäquat einschätzen zu können, wird das behandelnde Personal darum gebeten, bestimmte Informationen zu den Beratungs- und Behandlungsvorgängen zu erfassen.

Zu diesem Zweck wird ein Web-Link geteilt, welcher zu einem Online-Formular der Plattform KoboToolbox führt. Werden bei den Angeboten Patient:innen beraten oder behandelt, so kann vom Personal auf freiwilliger Basis dieses Formular ausgefüllt und abgesendet werden. Diese Datensätze werden anschließend von Seitens \ac{UVO} ausgewertet, um für bestimmte Treffen oder Versammlungen Statistiken zur aktuellen Lage geben zu können.

\subsection{Krankenhausumfeld}\label{sub:ukl}

\begin{figure}[h]
	\centering
	\includegraphics[width=\linewidth]{Images/cd_ukl}
	\caption[Kommunikationsdiagramm UKL]{Kommunikationsdiagramm zu einem Krankenhausumfeld am Beispiel des UKL}
	\label{fig:cdUKL}
\end{figure}

Im stationären und ambulanten Krankenhausumfeld gibt es bei dem Umgang mit wohnungslosen Personen zunächst keine grundlegenden Unterschiede im Vergleich zu Krankenkassenpatienten. Eine vorgeschaltete Instanz können dabei jedoch die Sozialarbeiter:innen sein, welche die jeweiligen Patient:innen mit ihren gesundheitlichen Anliegen an diese Gesundheitsinstitution verwiesen haben.

Demnach ist der Ablauf wie folgt: Der oder die Patient:in kommt zur Vorstellung und bringt dabei evtl. vorhandene Dokumente oder einen Überweisungsschein mit. Eine alternative Situation wäre im Rahmen einer Notfallvorstellung, das heißt die jeweilige Person kommt mittels des Rettungsdienstes oder wird zu späterer Uhrzeit vorstellig. In beiden Fällen wird wie gewohnt ein neuer Fall im Patientenverwaltungssystem angelegt, ggf. unter Verwendung von Pseudonymen oder sonstigen Platzhaltern.

Es erfolgt eine Behandlung und es wird mit dem oder der Patient:in abgesprochen, ob Folgetermine für eine Weiterbehandlung nötig oder gewünscht sind. Die physischen Dokumente der Patientenakte werden nach Abschluss eingescannt und in das digitale Archiv überführt. 

\subsection{VGP und SpDi}\label{sub:vgp}

\begin{figure}[h]
	\centering
	\includegraphics[width=\linewidth]{Images/cd_vgp}
	\caption[Kommunikationsdiagramm VGP]{Kommunikationsdiagramm für den Verbund Gemeindenahe Psychiatrie}
	\label{fig:cdVGP}
\end{figure}

Der \ac{VGP} umfasst die drei Angebote \acf{PIA}, \acf{SpDi} und \acf{TK}. Im Falle des \ac{SpDi} kann der erste Kontakt auf verschiedener Weise stattfinden: Das mobile Kontakt- und Beratungsteam erhält eine Meldung, meist vom Ordnungsamt oder von der Polizei, über eine Person auf der Straße, die sich psychisch auffällig verhält. Das Team macht sich auf den Weg und schätzt die Situation vor Ort ein, berät die gemeldete Person und vermittelt ggf. an eine passende Institution wie z.B. die \ac{PIA}. Anschließend erfolgt eine Rückmeldung über die durchgeführten Maßnahmen an die meldende Instanz. Andernfalls bietet der Dienst auch an fünf Standorten psychosoziale Sprechstunden an, zu denen die Personen selbstständig gehen können, wenn sie Hilfe benötigen.

Die Dokumentation erfolgt über \enquote{Octoware}, ein Anwendungssystem, welches auch beim Gesundheitsamt Einsatz findet, hierbei jedoch mit einem Modul speziell für \acp{SpDi} ausgestattet ist und sachsenweit von diesen benutzt wird. Die Beratungszeiten werden dabei als Episoden angelegt, in deren Lauf eine Dokumentation von soziodemografischen und Verwaltungsdaten sowie Diagnosen, Kontakten und evtl. vorhandenen Dokumenten erfolgt. Eine Behandlung erfolgt erst in der \ac{PIA}, wo die vorstellig werdenden Personen als eigener Fall angelegt und dokumentiert werden.

\subsection{Streetwork}\label{sub:sw}

\begin{figure}[h]
	\centering
	\includegraphics[width=\linewidth]{Images/cd_sw}
	\caption[Kommunikationsdiagramm Streetwork]{Kommunikationsdiagramm für Streetwork-Projekte}
	\label{fig:cdSW}
\end{figure}

Da sich Streetworker:innen primär um die allgemeinen Lebenssituationen von Wohnungs- und Obdachlosen kümmern, sind sie außerhalb von Behandlungen im Rahmen der ersten Hilfe, also Pflaster und Verbände für kleinere Verletzungen, auf den Rat und die Mitarbeit von medizinischem Personal angewiesen, wenn es um gesundheitliche Anliegen der Klient:innen geht. Eine Dokumentation im medizinischen Sinne erfolgt dabei nur sehr eingeschränkt im Austausch der Parteien untereinander.

Es bildet sich damit ein Dreieck bestehend aus den Informationsflüssen zwischen diesen drei Eckpunkten, in dem die Gruppe des ärztlichen Personals sich mit der medizinischen Versorgung der Betroffenen befasst und abhängig von der jeweiligen Situation - ehrenamtliche Angebote, in einer Praxis oder im Umfeld eines Krankenhauses - eigenständig dokumentiert. Der Austausch zwischen wohnungslosen Personen und den Streetworker:innen besteht neben den allgemeinen Unterstützungsangeboten auch darin, dass sie die Menschen auch zu Besuchen bei Ärzt:innen begleiten oder die relevanten Dokumente wieder an sie aushändigen.

In diesem Zusammenhang sprechen sich die Leistungserbringer untereinander ab, wobei es meist darum geht, dass sich eine der Parteien mit den Betroffenen in Verbindung setzt, um nach deren Wohlergehen zu schauen oder um Folgetermine für eine Behandlung zu vereinbaren. In anderer Sache informieren Streetworker:innen die ehrenamtliche Ärzt:innen auch, wenn sie jemanden mit gesundheitlichen Beschwerden an sie verwiesen haben.


\section{Modell des Informationssystems}\label{sec:3lgm2model}

Das nachfolgende Modell zielt darauf ab, das Informationssystem in Bezug auf die medizinische Dokumentation von Wohnungslosen als Ganzes darzustellen. In der fachlichen Ebene werden einige Aufgaben und Objekttypen auftauchen, die in ähnlicher Form bereits in den Kommunikationsdiagrammen zu finden sind. In diesem Fall wird jedoch verdeutlicht, bei welchen Aufgaben die entsprechenden Objekte entstehen bzw. welche Informationen dafür benötigt werden. Dies wird durch die zwei darunter liegenden Ebenen erweitert, indem gezeigt wird, welche Anwendungssysteme die einzelnen Elemente unterstützen oder tragen sowie auf welchen physischen Bausteinen diese wiederum bereitgestellt werden.

\subsection{Fachliche Ebene}

\begin{figure}[h]
	\centering
	\includegraphics[width=\linewidth]{Images/domain1}
	\caption[Fachliche Ebene - ehrenamtliche Versorgung]{Die fachliche Ebene der ehrenamtlichen Versorgung von Wohnungslosen}
	\label{fig:domain1}
\end{figure}

In der fachlichen Ebene (Abbildungen \ref{fig:domain1} und \ref{fig:domain2}) sind die bereits genannten Elemente so gruppiert, dass zu erkennen ist, welche Personen für diese zuständig sind. Eine essentielle Personengruppe sind die ehrenamtlichen Ärzt:innen, welche die Angebote des Projektes \ac{UVO} als Aufgabe haben. Der damit einhergehende Objekttyp \enquote{Kobo-Daten} beschreibt sämtliche Daten, die mittels der KoboToolbox erfasst werden. Die medizinisch relevanten Details sind diesem Element untergeordnet.

Ein weiterer Objekttyp, der in diesem Zusammenhang relevant ist, sind Pläne bzw. generell Informationen die vom Projekt \ac{UVO} ausgehen. Diese werden neben der Koordination der Angebote auch für die UVO-Treffen verwendet, in denen sich die einzelnen Interessenparteien austauschen können.

Die andere Aufgabe der Ärzt:innen ist die Anamnese während den CABL-Sprechstunden, wodurch die Fachrichtung des jeweiligen Falls geklärt wird. Die Kostenträgerermittlung wird ebenfalls für das Clearing benötigt. Der ausgehändigte \ac{ABS} kann anschließend bei der administrativen Aufnahme bei einer entsprechenden Gesundheitsinstitution verwendet werden.

\begin{figure}[h]
	\centering
	\includegraphics[width=\linewidth]{Images/domain2}
	\caption[Fachliche Ebene - UKL und SpDi]{Die fachliche Ebene mit dem Krankenhausumfeld und dem sozialpsychiatrischen Dienst}
	\label{fig:domain2}
\end{figure}

Die verschiedenen medizinisch relevanten Dokumente, die einige Patient:innen bei sich tragen und mitbringen, sind oft von Bedeutung, um sich einen Überblick über das Befinden der Person zu verschaffen. Unter Umständen kann es auch sein, dass die Dokumente von Streetworker:innen getragen werden, um sie von einem Arzt / einer Ärztin zu der jeweiligen Person oder andersherum zu bringen.

Die letzte Gruppe enthält die Aufgabe und die Objekttypen, die gegenüber dem \ac{SpDi} anfallen. Diese beziehen sich auf die beratenden und unterstützenden Angebote, die von dem Dienst ausgehen, und sind in ihrer Struktur ähnlich einer Fall- oder Patientenakte aufgebaut.

\subsection{Logische Werkzeugebene}

\begin{figure}[h]
	\centering
	\includegraphics[width=\linewidth]{Images/logical}
	\caption[Logische Werkzeugebene]{Die logische Werkzeugebene}
	\label{fig:logical}
\end{figure}

In dieser Werkzeugebene (Abbildung \ref{fig:logical}) ist zunächst auffällig, dass die Anzahl an Schnittstellen sehr begrenzt ist und die Verbindungen meist mit papierbasierten Anwendungsbausteinen realisiert werden.

Beispielsweise werden die Notizen, welche von ehrenamtlichen Ärzt:innen während den Behandlungen erstellt werden, zwar meist nur als Gedächtnisstütze genutzt, jedoch können diese auch per Transkription in ein Praxisverwaltungssystem überführt werden, um dort einen entsprechenden Fall anzulegen. Die ist zumindest der Fall, wenn versucht wird, die behandelten Personen an eine Arztpraxis anzubinden.

Im Falle von \ac{CABL} und \ac{UVO} sind die Bausteine in sich geschlossen und tauschen im Normalfall keine Daten miteinander aus. Über den \ac{ABS} als papierbasierter Baustein wird eine Verbindung mit dem Patientenverwaltungssystem, hier als Beispiel des \ac{UKL}, realisiert, wobei angemerkt werden muss, dass dieses weitere Schnittstellen zu anderen Anwendungssystemen besitzt, die für die Erfüllung von Aufgaben der Gesundheitsversorgung von Nöten sind. Dies wären beispielsweise ein Labor- oder ein Radiologieinformationssystem, die hierbei jedoch im Sinne der Anschaulichkeit weggelassen wurden.

Ein weiteres alleinstehendes Anwendungssystem wird zur Dokumentation der Arbeit des \acp{SpDi} verwendet und basiert auf der Software \enquote{Octoware}. Obwohl diese auch beim Gesundheitsamt eingesetzt und bei dieser Seite auch ein Interesse an einigen bestimmten Episoden oder sonstigen Daten besteht, existiert keine etablierte Schnittstelle zwischen diesen Instanzen.

\subsection{Physische Werkzeugebene}

\begin{figure}[h]
	\centering
	\includegraphics[width=\linewidth]{Images/physical}
	\caption[Physische Werkzeugebene]{Die physische Werkzeugebene}
	\label{fig:physical}
\end{figure}

Die unterste Ebene (Abbildung \ref{fig:physical}) beschreibt die physischen Bausteine, auf denen die Anwendungssysteme laufen, und deren Verbindungen zueinander. Auch hier ist zu erkennen, dass die einzelnen Komponenten ziemlich losgelöst voneinander arbeiten und Kommunikation hauptsächlich zwischen Clients, die auf eine Anwendung zugreifen wollen, und Servern / Rechenzentren, die diese bereitstellen, stattfindet.

Der humanitäre Server von KoboToolbox stellt beispielsweise die Plattform dar, auf die jene zugreifen können, die den entsprechenden Web-Link besitzen. In ähnlicher Weise wird die zur Ausstellung des \ac{ABS} benötigte Schein-Verwaltung auf den Servern des \ac{AKST} gehostet. In Bezug auf den \ac{VGP} oder einem Krankenhaus wie dem \ac{UKL}, werden die Anwendungen durch die eigene oder übergeordnete Institution bereitgestellt.


\section{Stärken und Schwächen}\label{sec:analysis}

Diese Herangehensweise an medizinische Angebote bildet in Bezug auf die dazugehörige Dokumentation ein heterogenes Bild, bestehend aus einer Vielzahl an Tools, die den jeweiligen Situationen entsprechend gewählt wurden und eingesetzt werden. Am Beispiel des \acp{SpDi} ist sogar eine Software im Einsatz, welche extra für die Dokumentation in diesem Aufgabenfeld vorgesehen ist.

Die Kommunikation zwischen den Leistungserbringern findet meist persönlich per Telefonat, E-Mail oder in Chatgruppen verschiedener Apps statt, wodurch ein Austausch sowohl zu medizinischen sowie zu sozialarbeitsbezogenen Themen auch kurzfristig stattfinden kann. Die Tools, welche im Zusammenhang mit der akuten Behandlung von Wohnungslosen eingesetzt werden, wie es z.B. bei den aufsuchenden Angeboten oft der Fall ist, wurden zudem so gewählt, dass dokumentarische Zeitaufwand für die behandelnden und beratenden Personen minimal ist.

Die Tatsache, dass einige Angebote (z.B. im Projekt \ac{UVO}) auf die Mithilfe von ehrenamtlich arbeitenden Ärzt:innen angewiesen sind, stellt gleichzeitig auch ein Nachteil in Bezug auf die medizinische Dokumentation dar. Da man die Ehrenamtler:innen nicht mehr belasten möchte als nötig, werden keine verpflichtenden Vorgaben zum Umfang und Inhalt gestellt. Somit ist es auch möglich, dass in einigen Fällen gar keine Informationen zu den Patient:innen oder den durchgeführten Leistungen vorliegen. Für spätere Behandlungen kann man sich anschließend nur auf Eigenaussagen der Betroffenen oder halt auf eigene Notizen stützen.

In Bezug auf die im Kapitel \ref{sec:criteria} definierten Kriterien, können auch weitere negative Eigenschaften des Leipziger Informationssystems erkannt werden. Im Modell \ref{fig:logical} ist zu erkennen, dass zwischen den einzelnen Anwendungssystemen so gut wie keine rechnerbasierten Schnittstellen existieren, über die Informationen ausgetauscht werden können, weshalb die Interoperabilität hierbei als niedrig eingestuft werden kann.

Zudem ist die Verfügbarkeit einzelner Anwendungen nicht immer gewährleistet, wie am Beispiel von Octoware, dem Dokumentationstool des \ac{SpDi}, welches nur an den festen Standorten verwendet werden kann. Auf Papier geschriebene Notizen werden in solchen Fällen genutzt, um diese zu einem späteren Zeitpunkt in das jeweilige Anwendungssystem zu überführen oder um den generellen Informationsfluss zwischen zwei solcher Systeme zu ermöglichen. In Bezug auf Kriterium K4, können hierbei relativ häufig Medienbrüche festgestellt werden.

Der andere Ansatzpunkt ist es, die Leute durch Behandlungsscheine, entweder durch das Sozialamt ausgestellt oder der \ac{ABS}, der von den Clearingstellen koordiniert wird, an die Regelversorgung anzubinden, was zu Folge hat, dass in den entsprechenden Institutionen eine Dokumentation, ggf. mit Pseudonym, von der Person hinterlegt wird. Aber auch hier gibt es den Nachteil, wenn eine pseudonymisierte Patientenakte bzw. die gespeicherten Daten eventuell nicht mehr einer vorstellig gewordenen Person zugeordnet werden können zumal die Scheine nur kurzfristig den Zugang zum Gesundheitssystem ermöglichen.


\section{Vergleich mit bestehenden Ansätzen}\label{sec:compare}

In Kapitel \ref{sec:approaches} wurden Lösungsansätze aufgeführt, von denen einige in Vergleich mit dem Leipziger Informationssystem zur medizinischen Dokumentation bei Wohnungs- und Obdachlosen gesetzt werden können. Die Gegenüberstellung mit der Wohlfahrt Safetynet und dem Hilfsprogramm BHCHP scheint hierfür besonders geeignet, da die Angebote beider Ansätze ebenfalls in der Größenordnung einer Stadt Anwendung finden. 

Mit Dublin als Sitz, stellt Safetynet einen zentralen Knotenpunkt bei der medizinischen Versorgung von Menschen dar, die sich außerhalb der Gesellschaft befinden. Das BHCHP nimmt für die Stadt Boston eine ähnliche Rolle ein, wobei dessen Ansatz auch die Verbesserung der Kooperation mit anderen Hilfsorganisationen vorsieht. Dem gegenüber besitzt Leipzig mehrere eigenständige Anlaufstellen, zu denen wohnungslose Personen mit gesundheitlichen Beschwerden gehen können, welche sich jedoch den jeweiligen Situationen entsprechend in ihren Anwendungssystemen und Verfahrensweisen unterscheiden können.

Eine zentrale Verwaltung hat auf einer übergeordneten fachlichen Ebene den Vorteil, dass ausgetauschte Daten einen gewissen Grad an Standardisierung haben. Im Falle von Safetynet wird ein webbasiertes Gesundheitsaktensystem eingesetzt, welches von sämtlichen Teams zur Dokumentation verwendet wird. Solange es keine Einschränkungen beim Betrieb des Tools gibt, kann sichergestellt werden, dass anfallende oder benötigte Daten an den relevanten Orten zur richtigen Zeit zur Verfügung stehen, ggf. auch über eine Schnittstelle mit den Anwendungssystemen eines Krankenhauses oder einer Arztpraxis.

BHCHP hat bei seinem Ansatz den Vorteil einer guten Interoperabilität. Da auf einen flächendeckenden Einsatz eines Patientenverwaltungssystems gesetzt wird, können die einzelnen Standorte bei entsprechender Anbindung frei und standardisiert miteinander kommunizieren. Dies führt ebenfalls zu einer Minimierung von Medienbrüchen, wodurch wiederum das Risiko für Übermittlungsfehler verringert wird.

Wie in den obigen Modellen aufgezeigt, sind im Leipziger Informationssystem eine Vielzahl alleinstehender Anwendungen zu finden, die darauf abzielen, den Bedürfnissen der Akteure möglichst gerecht zu werden. Zum Austausch mit den Anwendungssystemen anderer Gesundheitsinstitutionen wird häufig auf papierbasierte Bausteine wie handschriftliche Notizen, Behandlungsscheine oder Anamnesebögen zurückgegriffen, welche zwar eine gewisse Flexibilität aufweisen, aber dennoch für Leistungserbringer zu einem Mehraufwand führen, weil diese zunächst transkribiert werden müssen, damit sie in einem rechnerbasierten Anwendungssystem vor Ort verwendet werden können.

In Bezug auf die Qualität der Daten wäre sicherlich ein rechnerbasiertes Anwendungssystem mit entsprechenden Schnittstellen gut geeignet, um sowohl die Integrität und Authentizität als auch die generelle Korrektheit der Daten zu prüfen. Zwar gibt es auch bei papierbasierten Methoden Möglichkeiten, dies sicherzustellen, jedoch gehen sie auch mit einer erhöhten Belastung für die Leistungserbringer einher.

Mit Fokus auf die unterliegende logische Werkzeugebene zeigen sich beim Leipziger Informationssystem einige Vorteile gegenüber einem zentralisierten Ansatz, wie es bei Safetynet der Fall ist: Die Tatsache, dass kein Anwendungssystem allein für die Aufrechterhaltung des Informationsflusses zuständig ist, bedeutet auch, dass bei einem eventuellen Ausfall nur wenige Bereiche davon betroffen sind. Für solche Fälle sollte jedoch ohnehin ein entsprechendes Konzept vorliegen, um mit den einhergehenden Einschränkungen umzugehen.

Aus diesem Grund kann man auch hier davon ausgehen, dass die positiven Merkmale von Ansätzen, wie sie z.B. in Dublin eingesetzt werden - vor allem in Sachen Funktionalität und Nutzerfreundlichkeit - den erhöhten Risiken überwiegen, die mit einem zentralen Knotenpunkt für Kommunikation einhergehen, zumal es Vorkehrungen gibt, um diese zu minimieren.

Unter Betrachtung der physischen Ebene treten die Kriterien der Verfügbarkeit und Sicherheit mit in den Vordergrund. Diese sind im Leipziger Umfeld für die einzelnen Anwendungssysteme zu betrachten, sodass es einige positive Beispiele gibt, wie z.B. KoboToolbox, welche webbasiert abgerufen und somit auch über die privaten Mobilgeräte von Ärzt:innen vor Ort bedient werden kann. Die Daten werden dabei in einem für humanitäre Projekte dedizierten Rechenzentrum unter bestimmten Sicherheitsvorkehrungen gespeichert.

Jedoch ist dies nicht bei allen Tools der Fall. So wird häufig bei mobilen Einsätzen des \acp{SpDi} oder bei Angeboten, die durch ehrenamtliche Ärzt:innen betreut werden, auf handschriftliche Notizen zurückgegriffen, um diese zu einem späteren Zeitpunkt in das zuständige Anwendungssystem, z.B. Octoware oder ein Praxisverwaltungssystem, zu übertragen. Die Verwendung papierbasierter Elemente führt insbesondere im wohnungslosen Umfeld zu Bedenken hinsichtlich Datensicherheit und -schutz, vor allem, wenn obdachlose Personen für ihre medizinischen Dokumente zuständig sind, die in solch prekären Situationen schnell verloren gehen können.